% Glossaire des métriques et unités (type=metriques)
% Une entrée par clé ; les parents sont définis avant leurs enfants.
% Tri : métriques par nom, puis ressources FPGA (préfixe « zr- »), puis unités (préfixe « zz- »).

% ---------------------------------------------------------------- Précision de détection
\newglossaryentry{iou}{type=metriques, name={IoU}, sort={iou},
  description={Intersection over Union (indice de Jaccard) entre une boîte prédite $A$ et
  une boîte de référence $B$ : $\mathrm{IoU} = |A\cap B|/|A\cup B|$, de 0 à 1, plus haut =
  mieux. Sert de seuil pour décider qu'une détection est juste (AP, mAP). Employée seule
  comme précision, elle n'est pas comparable à une mAP : DAC-SDC note l'IoU moyenne d'un
  seul objet par image \citep[§2.3]{2021-xu-dac-sdc-low-power-object}, et SkyNet\pp{}
  passe de 0,741 à 0,727 puis 0,680 en quantifiant
  \citep[Table~7]{arxiv-1909.09709-skynet-hardware-efficient-method-object}. Un seuil de
  Jaccard sert aussi à définir une erreur critique de détecteur sous radiations
  \citep[§II-C]{2024-rech-artificial-neural-networks-space-safety}.}}

\newglossaryentry{ap}{type=metriques, name={AP}, sort={ap},
  description={Average Precision : aire sous la courbe précision-rappel d'une classe, les
  détections étant jugées justes au-delà d'un seuil d'IoU ; de 0 à 1 (ou en \%), plus haut =
  mieux. Le mode d'interpolation de la courbe varie selon le protocole (VOC, COCO)
  \horscorpus. Piège : dans la notation COCO, \enquote{AP} désigne la moyenne sur les
  classes et sur les seuils 0,50 à 0,95, c'est-à-dire la mAP50-95 ; une AP à une seule
  classe (véhicules) n'est pas comparable à une mAP multi-classes
  \citep[Table~8]{2023-zhai-fpga-vehicle-detection-tracking-accelerator}.}}

\newglossaryentry{ap50}{type=metriques, name={AP50}, sort={ap50}, parent=ap,
  description={AP au seuil d'IoU 0,5 (AP@0.5). Chez Zhai et al., l'AP@0.5 est identique en
  FP32 et en INT16 à trois décimales, ce qui laisse penser qu'elle n'a pas été remesurée
  sur la carte \citep[Table~8]{2023-zhai-fpga-vehicle-detection-tracking-accelerator}
  \deduit.}}

\newglossaryentry{map}{type=metriques, name={mAP}, sort={map},
  description={mean Average Precision : moyenne des AP par classe ; plus haut = mieux. Sans
  seuil précisé, le sens dépend du jeu (mAP50 sur VOC, mAP50-95 sur COCO). Pièges relevés :
  le jeu de données (COCO à 80 classes, VOC à 20, une classe filtrée de VOC chez LP-YOLO\pp{}
  \citep[§V]{arxiv-2609.24757-npu-accelerator-quantized-real-time}, jeu privé à une cible
  \citep[Table~4]{arxiv-2502.15737-performance-analysis-you-only-look}) ; la résolution
  d'entrée ; la mesure simulée et non faite sur la carte (quantification simulée par ONNX
  Runtime \citep[§IV-A]{2023-montgomerie-corcoran-satay-streaming-architecture-toolflow-accelerating}) ;
  seules des pertes relatives données \citep[§6.2]{2019-ding-req-yolo-resource-aware-efficient}.
  La mAP ne mesure pas la qualité de l'incertitude
  \citep[§IV-D]{arxiv-2011.10671-review-comparative-probabilistic-object-detection}.}}

\newglossaryentry{map50}{type=metriques, name={mAP50}, sort={map50}, parent=map,
  description={mAP au seuil d'IoU 0,5 (notée aussi mAP@0.5, mAP@0,5) ; métrique principale
  de VOC. Plus indulgente que la mAP50-95 sur la localisation : les deux ne se comparent pas.
  Exemple : 59,4 de mAP@0,5 sur les seuls véhicules de VOC pour LP-YOLO Slim\pp{}
  \citep[§V]{arxiv-2609.24757-npu-accelerator-quantized-real-time} ; 73,8 sur VOC, 1,5 point
  sous le flottant, sans dire si la mesure est faite sur la carte
  \citep[Table~2]{2024-yan-fpga-yolov5-accelerator-real-time}.}}

\newglossaryentry{map5095}{type=metriques, name={mAP50-95}, sort={map5095}, parent=map,
  description={mAP moyennée sur les seuils d'IoU de 0,50 à 0,95 par pas de 0,05 (notée aussi
  mAP@0.5:0.95) ; métrique principale de COCO. Mesurée sur la carte, elle révèle des pertes
  fortes et irrégulières : TensorRT INT8 38,6 contre DPU INT8 37,5 pour YOLOv4
  \citep[Table~1]{arxiv-2304.11580-framework-benchmarking-real-time-embedded} ; YOLOv8n
  de 0,64 à 0,34 sur KITTI, YOLOv11n de 0,75 à 0,71 sur un jeu de personnes
  \citep[Table~II]{arxiv-2607.13106-no-attention-no-problem-dpu}\pp ; 84,2 à 71,9 en INT8 sur
  Jetson \citep[Table~4]{arxiv-2502.15737-performance-analysis-you-only-look}.}}

\newglossaryentry{fppi}{type=metriques, name={FPPI}, sort={fppi},
  description={False Positives Per Image : nombre moyen de fausses détections par image.
  Le taux d'échec (miss rate) se lit à FPPI fixé ; le mMR du protocole Caltech en fait une
  moyenne logarithmique, plus bas = mieux. Sur CrowdHuman, un FPN atteint un mMR de 50,42
  sur le corps entier contre 10,08 sur Caltech
  \citep[Table~7]{arxiv-1805.00123-crowdhuman-benchmark-detecting-human-crowd}\pp. Le mMR
  reste une moyenne : il ne dit pas le taux de personnes manquées à un instant ni le pire
  cas, ce qui compte pour la sûreté \deduit.}}

% ---------------------------------------------------------------- Classification, anomalies
\newglossaryentry{auroc}{type=metriques, name={AUROC}, sort={auroc},
  description={Aire sous la courbe ROC (taux de vrais positifs contre taux de faux
  positifs), notée aussi AUC ; de 0,5 (hasard) à 1, souvent donnée en \%. Indépendante du
  seuil, elle ne dit donc rien du point de fonctionnement retenu ; le survey de Cui et al.
  ne traite pas le réglage du seuil \citep[§VI-B]{2023-cui-survey-unsupervised-anomaly-detection-algorithms}
  (AUROC image de 99,1 pour PatchCore sur MVTec AD, 92,2 sur un jeu réel, Tables~10 et 12).
  Piège : AUC évaluées en TensorFlow et non sur le cœur FPGA
  \citep[§5--6]{2019-tsukada-neural-network-device-learning-anomaly}.}}

\newglossaryentry{f1}{type=metriques, name={F1}, sort={f1},
  description={Moyenne harmonique de la précision $P$ et du rappel $R$ :
  $F_1 = 2PR/(P+R)$, de 0 à 1 (ou en \%), plus haut = mieux ; dépend du seuil de décision.
  Exemple de perte à la quantification : F1 nul après PTQ INT8 seule, 92,77 après
  fine-tuning contre 94,05 en FP32, et 95,29 sur Xavier NX
  \citep[Tables~III, VI]{arxiv-2407.17647-energy-efficient-artefact-detection-accelerator}.}}

% ---------------------------------------------------------------- Compression d'images
\newglossaryentry{psnr}{type=metriques, name={PSNR}, sort={psnr},
  description={Peak Signal-to-Noise Ratio : $10\log_{10}(\mathrm{MAX}^2/\mathrm{MSE})$, en dB,
  plus haut = mieux. Ne se lit qu'avec le débit (bpp) : un point débit-distorsion, pas une
  qualité absolue. Exemple : 29,94\,dB à 0,2542\,bpp pour F-LIC sur VCU118
  \citep[Fig.~6]{2022-sun-f-lic-fpga-learned-image}.}}

\newglossaryentry{ssim}{type=metriques, name={SSIM}, sort={ssim},
  description={Structural Similarity : indice de similarité perceptuelle (luminance,
  contraste, structure), au plus 1, plus haut = mieux ; MS-SSIM en est la version
  multi-échelle. Exemple : 0,9647 sur VPU Myriad~2 contre 0,9651 en référence
  \citep[§III-B--C]{2023-guerrisi-artificial-intelligence-board-image-compression}.}}

\newglossaryentry{bpp}{type=metriques, name={bpp}, sort={bpp},
  description={Bits par pixel du flux compressé ; plus bas = meilleure compression à
  qualité égale. Piège : le débit n'a de sens que si le codage entropique est compté ; chez
  F-LIC, il tourne sur le PC hôte \citep[Fig.~1]{2022-sun-f-lic-fpga-learned-image}, et
  Φ-Sat-2 compresse d'un facteur 8 fixe, sans codage entropique
  \citep[§II-A]{2023-guerrisi-artificial-intelligence-board-image-compression}.}}

\newglossaryentry{bd-rate}{type=metriques, name={BD-rate}, sort={bd-rate},
  description={Bjøntegaard delta rate : écart moyen de débit, en \%, entre deux courbes
  débit-distorsion à qualité égale ; négatif = gain, positif = perte. La perte d'un
  Ballé2018 quantifié en INT8 passe de 30,0 à 6,30\,\%, en simulation seulement
  \citep[Table~I]{arxiv-2511.17135-multi-stage-optimization-framework-deploying}\pp.}}

% ---------------------------------------------------------------- Débit et efficacité
\newglossaryentry{fps}{type=metriques, name={FPS}, sort={fps},
  description={Images (frames) traitées par seconde ; plus haut = mieux. Ce n'est pas
  l'inverse de la latence dès que le calcul est pipeliné ou groupé (une \enquote{latence}
  de 5,6\,ns qui est en fait un débit amorti, pour 85\,ns réels
  \citep[§IV-E]{2024-que-low-latency-variational-autoencoder-fpgas}). Pièges : FPS du cœur
  du réseau ou de bout en bout (NMS et pré-traitement sur l'ARM exclus
  \citep[§2.5]{arxiv-2211.11647-benchmarking-edge-computing-devices-grape}, post-traitement
  mono-cœur qui limite à 34,05\,FPS \citep[§IV-C]{arxiv-2607.13106-no-attention-no-problem-dpu}\pp) ;
  batch (batch~4 sur FPGA pour SkyNet\pp{}
  \citep[Tables~5--7]{arxiv-1909.09709-skynet-hardware-efficient-method-object}, batch~1
  contre 8 \citep[Table~VI]{2024-aggarwal-chameleon-dual-memory-replay-online}) ;
  résolution d'entrée et jeu (256--320 pixels contre 640).}}

\newglossaryentry{fpsw}{type=metriques, name={FPS/W}, sort={fps/w}, parent=fps,
  description={Débit par watt, $\mathrm{FPS}/P$, soit des images par joule ; son inverse est
  l'énergie par image. Plus haut = mieux. N'a de sens qu'à portée de puissance égale
  (colonne \texttt{power\_scope} : puce, carte, carte moins le repos, estimée). À même
  méthode (carte au multimètre), YOLOv4 INT8 tient 1,31\,FPS/W sur ZCU104 contre 2,05 sur AGX
  Xavier \citep[Table~1]{arxiv-2304.11580-framework-benchmarking-real-time-embedded} ;
  2,6\,FPS/W présentés comme efficaces sont sous quatre des cinq travaux cités
  \citep[§IV-B, Table~I]{arxiv-2507.18174-real-time-object-detection-classification}.}}

\newglossaryentry{gop}{type=metriques, name={GOP}, sort={gop},
  description={Giga-opérations : nombre d'opérations (10$^9$) d'une inférence, propriété du
  modèle et de la résolution, pas du matériel ; à ne pas confondre avec GOPS (un débit).
  Exemple : 5,7\,GOP pour YOLOv3-tiny, en comptant un MAC pour deux opérations
  \citep[Table~5]{2024-zhang-end-end-acceleration-yolo-object}.}}

\newglossaryentry{gops}{type=metriques, name={GOPS}, sort={gops},
  description={Giga-opérations par seconde (GOP/s), débit de calcul ; plus haut = mieux. Pour
  un réseau, $\mathrm{GOPS} = \mathrm{GOP} \times \mathrm{FPS}$. Pièges : 1 MAC compte pour
  2 opérations, convention pas toujours rappelée
  \citep[Table~5]{2024-zhang-end-end-acceleration-yolo-object} ; GOPS de crête contre GOPS
  effectifs sur un réseau ; GOPS projetés sur une FPGA hypothétique comparés à des GOPS mesurés
  \citep[Tables~16--17]{2019-coutinho-deep-neural-network-hardware-implementation} ; GOPS
  sans FPS pour un détecteur \citep[Table~5]{arxiv-1809.04570-finn-r-end-end-deep}, ou
  \enquote{10,65x} entre réseaux de tailles différentes à FPS voisins
  \citep[§V-D]{arxiv-2403.18921-smof-streaming-modern-cnns-fpgas}\pp.}}

\newglossaryentry{gopsw}{type=metriques, name={GOPS/W}, sort={gops/w}, parent=gops,
  description={Efficacité énergétique de calcul ; plus haut = mieux. Hérite des pièges des
  GOPS (crête ou effectifs) et de ceux de la puissance (portée) : GOPS/W calculés sur une
  puissance dynamique comparés à des GOPS/W totaux
  \citep[Table~8]{2023-zhai-fpga-vehicle-detection-tracking-accelerator}.}}

\newglossaryentry{gopsdsp}{type=metriques, name={GOPS/DSP}, sort={gops/dsp}, parent=gops,
  description={Débit rapporté au nombre de DSP utilisés, censé mesurer le rendement de
  l'architecture ; plus haut = mieux. Valable seulement dans une même famille de FPGA
  \citep[§3.1]{2018-venieris-toolflows-mapping-convolutional-neural-networks}. Elle efface
  la précision : 4 multiplications 4~bits par DSP donnent 0,49\,GOPS/DSP
  \citep[Table~3]{2024-yan-fpga-yolov5-accelerator-real-time}, un MAC INT16 par DSP
  0,018 \citep[Table~9]{2025-kim-design-implementation-yolov2-accelerator-zynq} ;
  comparaisons de W1A8 à W32A32
  \citep[Table~III]{2023-montgomerie-corcoran-satay-streaming-architecture-toolflow-accelerating}.}}

\newglossaryentry{tops}{type=metriques, name={TOPS}, sort={tops},
  description={Téra-opérations par seconde (10$^{12}$ op/s). Dans le corpus, presque
  toujours une valeur de crête : 275\,TOPS annoncés par le constructeur pour une Jetson AGX
  Orin \citep[Table~4]{arxiv-2411.01431-efficient-deep-learning-infrastructures-embedded} ;
  66\,TOPS de pic binaire sur ZU19EG, environ 16 fois le pic 8~bits
  \citep[§3.1]{arxiv-1612.07119-finn-framework-fast-scalable-binarized}. Une crête dépend
  de la précision et ne prédit pas le débit sur un réseau réel \deduit.}}

\newglossaryentry{flops}{type=metriques, name={FLOPs}, sort={flops},
  description={Nombre d'opérations en virgule flottante d'un modèle (un compte, pas un
  débit). L'usage confond souvent FLOPs et MAC (facteur 2) \horscorpus. Les FLOPs ne
  prédisent pas la latence réelle
  \citep[§2.3]{arxiv-2411.01431-efficient-deep-learning-infrastructures-embedded}.}}

\newglossaryentry{gflops}{type=metriques, name={GFLOPS}, sort={gflops},
  description={Giga-opérations flottantes par seconde : le GOPS d'un calcul en virgule
  flottante. Exemples : 61,62\,GFLOPS en FP32 sur Virtex-7
  \citep[Tables~7--8]{2015-zhang-optimizing-fpga-accelerator-design-deep} ; 28,15\,GFLOPS
  en apprentissage sur ZCU102, à batch~128
  \citep[§6]{arxiv-2202.10935-ef-train-enable-efficient-device}\pp.}}

\newglossaryentry{mac}{type=metriques, name={MAC}, sort={mac},
  description={Multiply-Accumulate : une multiplication suivie d'une addition, opération de
  base des convolutions ; compte pour 2 opérations dans les GOPS
  \citep[Table~5]{2024-zhang-end-end-acceleration-yolo-object}. Le nombre de MAC par DSP
  dépend de la précision : un MAC INT16
  \citep[§3.1]{2025-kim-design-implementation-yolov2-accelerator-zynq}, deux en 8~bits
  \citep[§3.1]{2024-zhang-end-end-acceleration-yolo-object}, quatre en 4~bits
  \citep[§4--5]{2024-yan-fpga-yolov5-accelerator-real-time}.}}

% ---------------------------------------------------------------- Latence et temps réel
\newglossaryentry{p99}{type=metriques, name={p99}, sort={p99},
  description={99\textsuperscript{e} centile d'une distribution (ici de latence) : 99\,\% des
  mesures sont en dessous. Plus informatif qu'une moyenne pour le temps réel, sans être un
  pire cas. Aucun papier du corpus ne donne de p99, de pire cas ni de gigue de bout en bout :
  moyennes sur quelques passes ou une vidéo
  \citep{arxiv-2304.11580-framework-benchmarking-real-time-embedded}, latence variable selon
  la vidéo \citep[Fig.~14]{2026-fata-low-cost-fpga-enhanced-cnn}.}}

\newglossaryentry{wcet}{type=metriques, name={WCET}, sort={wcet},
  description={Worst-Case Execution Time : borne supérieure, obtenue par analyse statique,
  du temps d'exécution d'un programme sur une cible donnée ; exigé par DO-178C. Dans le
  corpus, seul ACETONE le calcule (OTAWA), sur un lpc2138 non représentatif de la cible
  mesurée, et l'analyse échoue pour AlexNet
  \citep[§4.2, §5]{2022-de-albuquerque-silva-acetone-predictable-programming-framework-ml}.
  Le WCET entre PS et PL n'est pas traité
  \citep[§III-B]{arxiv-2509.04153-real-time-fpga-cnns-detection}.}}

% ---------------------------------------------------------------- Sûreté et certification
\newglossaryentry{fit}{type=metriques, name={FIT}, sort={fit},
  description={Failures In Time : nombre de défaillances pour $10^9$ heures de
  fonctionnement ; plus bas = mieux. Li et al. calculent
  $\mathrm{FIT} = R_\mathrm{raw} \times \text{taille} \times P(\mathrm{SDC})$ avec
  $R_\mathrm{raw} = 20{,}49$\,FIT/Mb extrapolé à 16\,nm \citep[§4.7]{2017-li-understanding-error-propagation-deep-learning},
  citent la cible de moins de 10\,FIT d'ISO~26262 \citep[§1]{2017-li-understanding-error-propagation-deep-learning}
  et ramènent le FIT d'Eyeriss de 8,55 à 0,35 \citep[§6.2]{2017-li-understanding-error-propagation-deep-learning}.
  Le passage d'un FIT mesuré à un SIL ou un PL n'est traité nulle part \deduit.}}

\newglossaryentry{sil}{type=metriques, name={SIL}, sort={sil},
  description={Safety Integrity Level (IEC 61508) : niveau d'intégrité de sécurité d'une
  fonction, de SIL~1 à SIL~4 (le plus exigeant), associé à une probabilité de défaillance
  dangereuse par heure ou à la sollicitation. Dans le corpus, l'IEC 61508 n'apparaît que
  comme contexte \citep[§1, §2.3]{2017-li-understanding-error-propagation-deep-learning} ;
  l'équivalent automobile (ASIL, ISO~26262) est plus présent
  \citep[§III-C]{2024-rech-artificial-neural-networks-space-safety}, et aucune fiche ne
  relie une précision de détection ou un FIT à un SIL \deduit.}}

\newglossaryentry{dal}{type=metriques, name={DAL}, sort={dal},
  description={Design Assurance Level : niveau d'assurance de conception en aéronautique,
  de A (défaillance catastrophique) à E, défini par DO-178C pour le logiciel et DO-254 pour
  le matériel électronique complexe (FPGA, ASIC). Dans le corpus : la conformité DO-254 des
  FPGA est hors périmètre de CoDANN
  \citep[§4.2.3]{2020-easa-concepts-design-assurance-neural-networks} ; CoDANN~II renvoie à
  DO-178C §12.2 pour la qualification des outils
  \citep[§3.7.1]{2024-easa-concepts-design-assurance-neural-networks} ; l'EASA renvoie les
  FPGA à AMC 20-152A \citep[MOC IMP-05]{2024-easa-concept-paper-guidance-level-1}.}}

% ---------------------------------------------------------------- Protocole
\newglossaryentry{dac-sdc}{type=metriques, name={DAC-SDC}, sort={dac-sdc},
  description={System Design Contest de la conférence DAC : détection d'un seul objet par
  image (images de drone), sur PYNQ-Z1 et Jetson TX2, mesurée par les organisateurs. Score
  $= \mathrm{IoU} \times (1 + \mathrm{ES})$, ES logarithmique en énergie avec des bases
  différentes pour FPGA et GPU ; débit minimal de 5\,FPS sur FPGA et 20 sur GPU ; FPGA au
  wattmètre, TX2 par ses capteurs internes \citep[§2.3]{2021-xu-dac-sdc-low-power-object}.
  Son IoU ne se compare pas à une mAP, et le classement ne passe pas d'une catégorie à
  l'autre.}}

% ---------------------------------------------------------------- Ressources FPGA
\newglossaryentry{lut}{type=metriques, name={LUT}, sort={zr-lut},
  description={Look-Up Table : cellule logique programmable de la FPGA (6 entrées sur les
  familles AMD récentes \horscorpus), qui réalise la logique et de petites mémoires. Les
  architectures streaming déplacent le calcul des DSP vers les LUT : l'accumulation par
  popcount coûte à peu près moitié moins de LUT qu'une accumulation signée
  \citep[§4.2.1]{arxiv-1612.07119-finn-framework-fast-scalable-binarized}.}}

\newglossaryentry{ff}{type=metriques, name={FF}, sort={zr-ff},
  description={Flip-flop : bascule de registre de la logique programmable. Un gain de FF
  n'a de sens qu'à débit égal : \enquote{$-26$\,\% de FF} face à des accélérateurs de 10,5 à
  62,9\,GOPS, pour 2,33\,GOPS
  \citep[Table~9]{2025-kim-design-implementation-yolov2-accelerator-zynq}.}}

\newglossaryentry{dsp}{type=metriques, name={DSP}, sort={zr-dsp},
  description={Bloc arithmétique câblé (DSP48 chez AMD : multiplieur et accumulateur), la
  ressource rare des petites FPGA (220 sur le XC7Z020 de la PYNQ-Z1
  \citep[§2.2]{2021-xu-dac-sdc-low-power-object}). La quantification change son
  rendement (voir GOPS/DSP) ; la QAT et l'élagage font passer un modèle de 6\,377 à 174 DSP
  \citep[Table~3]{arxiv-2101.05108-fast-convolutional-neural-networks-fpgas}.}}

\newglossaryentry{bram}{type=metriques, name={BRAM}, sort={zr-bram},
  description={Block RAM : mémoire on-chip en blocs de 36\,Kb, divisibles en deux de 18\,Kb
  \horscorpus ; comptée en blocs dans \texttt{benchmarks.csv}, pas en Mo. Piège : un nombre de
  BRAM dit peu de la capacité utilisée, la forme des mémoires de poids s'accordant mal aux
  blocs (52,9\,\% d'occupation pour un ResNet-50
  \citep[Table~IV]{arxiv-2011.07317-memory-efficient-dataflow-inference-deep}\pp). Les BRAM
  sont exclues de l'injection de fautes de Libano et al.
  \citep[§III-A]{2019-libano-selective-hardening-neural-networks-fpgas}.}}

\newglossaryentry{uram}{type=metriques, name={URAM}, sort={zr-uram}, parent=bram,
  description={UltraRAM : blocs de mémoire on-chip plus grands (288\,Kb) des
  UltraScale+ \horscorpus. Sur KV260, un design place les couches 11 à 21 en URAM et occupe
  100\,\% des BRAM et des URAM \citep[Table~9]{2026-fata-low-cost-fpga-enhanced-cnn}.}}

% ---------------------------------------------------------------- Unités
\newglossaryentry{u-w}{type=metriques, name={W}, sort={zz-01-w},
  description={Watt, unité de puissance. Dans le corpus, de quelques mW (puissance dynamique
  \citep[Table~16]{2019-coutinho-deep-neural-network-hardware-implementation}) et 0,04\,W
  sur MCU GAP9 \citep[Table~II]{arxiv-2403.07851-12-mj-per-class-device} à 2--23\,W sur SoC
  FPGA et environ 120--134\,W pour deux Alveo U280
  \citep[Table~4]{2021-alonso-elastic-df-scaling-performance-dnn}. Toujours lire la portée :
  une ZC706 consomme environ 7\,W au repos, inclus dans une mesure de carte
  \citep[§5.2.2]{arxiv-1612.07119-finn-framework-fast-scalable-binarized}.}}

\newglossaryentry{u-mw}{type=metriques, name={mW}, sort={zz-02-mw},
  description={Milliwatt ($10^{-3}$\,W), ordre de grandeur des MCU et ASIC : 62\,mW mesurés
  sur la puce VEGA \citep[Table~IV]{arxiv-2110.10486-tinyml-platform-device-continual-learning}.}}

\newglossaryentry{u-j}{type=metriques, name={J}, sort={zz-03-j},
  description={Joule, unité d'énergie : $E = P \times t = P/\mathrm{FPS}$ par image ; plus
  bas = mieux. Exemples : 1,41\,J par image 720p sur VCU118
  \citep[Fig.~6]{2022-sun-f-lic-fpga-learned-image} ; 1,22\,J par image d'apprentissage sur
  ZCU102 \citep[Table~VI]{2024-aggarwal-chameleon-dual-memory-replay-online}.}}

\newglossaryentry{u-mj}{type=metriques, name={mJ}, sort={zz-04-mj},
  description={Millijoule. Exemples : 21,52\,mJ par inférence sur DPU ZCU104
  \citep[Table~VI]{arxiv-2407.17647-energy-efficient-artefact-detection-accelerator} ;
  11 à 23\,mJ par nouvelle classe sur MCU \citep[Tables~II, IV]{arxiv-2403.07851-12-mj-per-class-device}.}}

\newglossaryentry{u-uj}{type=metriques, name={µJ}, sort={zz-05-uj},
  description={Microjoule ($10^{-6}$\,J), énergie d'une inférence de petit réseau :
  196\,µJ sur MAX78000 \citep[§III]{arxiv-2306.00001-tinyissimoyolo-quantized-low-memory-footprint} ;
  2\,535 et 44\,330\,µJ pour FINN et hls4ml sur Pynq-Z2
  \citep[Table~5]{arxiv-2206.11791-open-source-fpga-ml-codesign}\pp.}}

\newglossaryentry{u-s}{type=metriques, name={s}, sort={zz-06-s},
  description={Seconde. Rencontrée pour les traitements lents : environ 12,6\,s par image
  pour YOLOv2 sur XC7Z020 \citep[§4.2]{2025-kim-design-implementation-yolov2-accelerator-zynq} ;
  2,07\,s par événement d'apprentissage sur VEGA
  \citep[Table~IV]{arxiv-2110.10486-tinyml-platform-device-continual-learning}.}}

\newglossaryentry{u-ms}{type=metriques, name={ms}, sort={zz-07-ms},
  description={Milliseconde, unité usuelle de latence d'un détecteur : de 15,2\,ms de bout
  en bout \citep[Table~5]{2024-zhang-end-end-acceleration-yolo-object} à 54,42\,ms
  \citep[Table~1]{arxiv-2304.11580-framework-benchmarking-real-time-embedded}. Préciser
  ce que couvre la latence (inférence seule, PL seule, bout en bout).}}

\newglossaryentry{u-us}{type=metriques, name={µs}, sort={zz-08-us},
  description={Microseconde : latence d'un petit réseau hls4ml, environ 5\,µs sur VU9P et
  171\,µs sur XC7Z020 \citep[Tables~3--4]{arxiv-2101.05108-fast-convolutional-neural-networks-fpgas}.}}

\newglossaryentry{u-ns}{type=metriques, name={ns}, sort={zz-09-ns},
  description={Nanoseconde : latences de réseaux minuscules en logique pure, 85\,ns
  (17 cycles) pour un VAE dense \citep[§IV-E]{2024-que-low-latency-variational-autoencoder-fpgas},
  233\,ns chez Libano et al.
  \citep[§V]{2020-libano-understanding-impact-quantization-accuracy-radiation}.}}

\newglossaryentry{u-mhz}{type=metriques, name={MHz}, sort={zz-10-mhz},
  description={Mégahertz, fréquence d'horloge ; débit et puissance en dépendent. Dans le
  corpus, de 100 à 300\,MHz pour la logique FPGA ; un Tegra~X1 bridé à 76,8\,MHz pour
  égaliser le budget de puissance
  \citep[§3.1]{arxiv-1711.08740-fpgaconvnet-toolflow-mapping-diverse-convolutional}.}}

\newglossaryentry{u-ko}{type=metriques, name={Ko}, sort={zz-11-ko},
  description={Kilo-octet (KB, kB) ; le corpus ne distingue pas 1\,000 et 1\,024 octets.
  Ordre de grandeur des SRAM de MCU : entraînement sous 256\,KB, 206\,KB mesurés
  \citep[§3.2]{arxiv-2206.15472-device-training-256kb-memory} ; 630\,Ko de BRAM sur
  PYNQ-Z1 \citep[§2.2]{2021-xu-dac-sdc-low-power-object}.}}

\newglossaryentry{u-mo}{type=metriques, name={Mo}, sort={zz-12-mo},
  description={Mégaoctet (MB) ; 1\,Mo = 8\,Mbit, à ne pas confondre avec le Mbit des fiches
  techniques. Exemple : 3,5\,Mo de tampons de rejeu contre 48\,Mo
  \citep[Tables~II, VI]{2024-aggarwal-chameleon-dual-memory-replay-online}.}}

\newglossaryentry{u-mbit}{type=metriques, name={Mbit}, sort={zz-13-mbit},
  description={Mégabit, unité des capacités mémoire on-chip des FPGA. Un YOLOv3-tiny
  demanderait environ 34,7\,Mbit de poids à 4~bits, pour environ 4,9\,Mbit de BRAM sur le
  XC7Z020 \citep[§III-A]{arxiv-2609.24757-npu-accelerator-quantized-real-time}\pp.}}
